\chapter{Conclusão}

O projeto Civitas foi desenvolvido a partir da necessidade de organizar informações administrativas e financeiras da gestão pública municipal, considerando as dificuldades associadas à utilização de registros dispersos, à falta de integração entre cadastros e às limitações no acompanhamento de despesas e orçamentos. Nesse contexto, estabeleceu-se como objetivo a construção de uma solução web voltada à centralização dessas informações, ao apoio às rotinas administrativas e à rastreabilidade das operações. A proposta relaciona a organização dos dados às necessidades de planejamento, consulta e controle dos recursos públicos.

Ao longo do trabalho, o levantamento de requisitos permitiu delimitar o escopo da solução e identificar os serviços previstos para os diferentes perfis de acesso. A modelagem do sistema, por sua vez, organizou as entidades, os relacionamentos e os comportamentos necessários à representação do domínio administrativo e financeiro. A articulação entre requisitos, casos de uso e dicionário de classes contribuiu para explicitar as responsabilidades dos usuários e os vínculos entre os registros, oferecendo uma base para orientar o desenvolvimento e a manutenção do Civitas.

Entre os resultados documentados, destacam-se a definição dos requisitos funcionais e não funcionais, a elaboração da modelagem UML e a disponibilização dos scripts da estrutura relacional do banco de dados. Esses artefatos representam entregas concretas de especificação e estruturação da solução. A organização dos dados estabelece a relação entre secretarias, instituições, fornecedores, orçamentos, unidades consumidoras e despesas, além de contemplar informações de usuários e auditoria. Dessa forma, o trabalho avançou na representação integrada dos elementos necessários ao controle administrativo e financeiro, sem que a existência desses artefatos implique, isoladamente, a conclusão de todas as funcionalidades previstas.

A estruturação do domínio também permitiu esclarecer aspectos relevantes para a consistência do sistema. O planejamento orçamentário foi relacionado às instituições e aos tipos de despesa, enquanto as unidades consumidoras passaram a representar o contexto administrativo e financeiro dos lançamentos. Essa organização oferece uma base para associar despesas aos respectivos responsáveis e às informações necessárias ao seu acompanhamento. Trata-se de uma contribuição técnica para a construção da solução, cuja efetividade operacional depende da correspondência entre os modelos definidos e os componentes efetivamente implementados.

Os benefícios pretendidos pelo Civitas incluem maior facilidade de consulta, redução da dependência de controles fragmentados e melhores condições para acompanhar informações orçamentárias e financeiras. Entretanto, tais benefícios devem ser compreendidos como resultados esperados de sua utilização, e não como impactos já comprovados. A documentação disponível não apresenta evidências suficientes para afirmar redução mensurada de custos, melhoria de desempenho administrativo ou validação integral por usuários em ambiente real. Assim, a avaliação do projeto deve distinguir as entregas de especificação e estruturação já documentadas dos efeitos que dependem de implantação e acompanhamento.

No âmbito acadêmico, o trabalho possibilitou relacionar conhecimentos de análise de requisitos, modelagem de sistemas e organização de bancos de dados a um problema específico da administração municipal. A definição de entidades, regras e associações exigiu considerar a correspondência entre necessidades administrativas e representações técnicas, evidenciando que o desenvolvimento de uma solução não se limita à criação de funcionalidades isoladas. A coerência entre os artefatos e a clareza na delimitação das responsabilidades constituíram aspectos importantes da elaboração do projeto.

A análise do estágio documentado também evidencia a necessidade de continuidade na consolidação dos componentes e na verificação de sua correspondência com os requisitos. A comprovação do atendimento às metas de segurança, desempenho, acessibilidade e usabilidade depende de avaliações compatíveis com os critérios estabelecidos para o Civitas. Da mesma forma, a conclusão das funcionalidades e sua adequação às rotinas municipais precisam ser demonstradas por evidências de implementação e validação. Essas condições delimitam o alcance dos resultados apresentados e orientam a continuidade do trabalho sem antecipar conclusões sobre sua utilização prática.

Conclui-se que o Civitas apresenta uma proposta estruturada de apoio à gestão pública municipal, com resultados documentados na definição de requisitos, na modelagem do domínio e na organização dos dados administrativos e financeiros. Sua relevância está na aproximação entre o conhecimento acadêmico e a necessidade de integrar informações sobre instituições, orçamentos e despesas em uma mesma solução. A continuidade do desenvolvimento e a verificação dos resultados operacionais deverão permitir avaliar em que medida essa base técnica se traduz em apoio efetivo ao planejamento e ao controle dos recursos públicos municipais.
